  \subsubsection{整数への丸めで不動点と停止を調べる}\label{節：ガウス型の不動点と停止}
    丸めを繰り返す数列では,一般項だけでなく,どの整数で止まるかにも注目しよう.
    次の値が今の値と等しくなる値を不動点という.
    不動点に到達すると,その後は同じ値が続く.
    ただし,丸めを外した実数の漸化式と,不動点が同じとは限らない.

    整数$K$を引いて,
    \[
      b_{n+1}=\left\lfloor\frac{b_n}{2}\right\rfloor
    \]
    に帰着できる場合がある.
    正整数$m$に対して,
    \[
      \left\lfloor\frac{\lfloor x\rfloor}{m}\right\rfloor
      =\left\lfloor\frac xm\right\rfloor
    \]
    が成り立つので,繰り返すと分母をまとめられる.
    この性質を確かめよう.\,$k=\lfloor x\rfloor$を$m$で割り,\,$k=mq+r$と書くと,
    \,$0\le r\le m-1$であり,\,$x=mq+r+t\ (0\le t<1)$となる.
    これより$0\le r+t<m$だから,どちらの床関数も$q$になる.

    \textbf{初項の位置と,負の数の床関数の値}も重要である.
    例えば,\,$-1<x<0$なら$\lfloor x\rfloor=-1$であって0ではない.

    \noindent \bookblocklink{example-gauss-role-6}{problem}{answer}{【例題】}
    \begin{enumerate}[label=(\arabic*),parsep=3pt,beginpenalty=10000,format=\bookitemlink{example-gauss-role-6}{problem}{answer}]
      \item 非負整数$M$に対し,数列を次のように定める.
            \[
              a_1=M,\qquad
              a_{n+1}=\left\lfloor\frac{a_n+6}{2}\right\rfloor
              \quad(n\ge1).
            \]
            一般項と,最終的に一定になる値を求めよ.
    \end{enumerate}

    \noindent \bookblocklink{example-gauss-role-6}{hint}{problem}{【方針】}
    \begin{enumerate}[label=(\arabic*),parsep=3pt,beginpenalty=10000,format=\bookitemlink{example-gauss-role-6}{hint}{problem}]
      \item 6を引くと,床関数によって半分にする漸化式になります.
            初項が6より小さい場合は,負の数の床関数を使って考えましょう.
    \end{enumerate}

    \noindent \bookblocklink{example-gauss-role-6}{answer}{problem}{【解答】}
    \begin{enumerate}[label=(\arabic*),parsep=3pt,beginpenalty=10000,format=\bookitemlink{example-gauss-role-6}{answer}{problem}]
      \item \[
              a_n=6+\left\lfloor\frac{M-6}{2^{n-1}}\right\rfloor.
            \]
            有限回の操作の後で,\,$M<6$なら5,\,$M\ge6$なら6に固定される.
    \end{enumerate}

    \noindent \bookblocklink{example-gauss-role-6}{solution}{problem}{【解法】}
    \begin{enumerate}[label=(\arabic*),parsep=3pt,beginpenalty=10000,format=\bookitemlink{example-gauss-role-6}{solution}{problem}]
      \item 初項からすべての項が非負整数となる.
            \,$b_n=a_n-6$とおくと,整数6を床関数の外に出せるので,
            \begin{align*}
              b_{n+1}
                &=\left\lfloor\frac{b_n+12}{2}\right\rfloor-6\\
                &=\left\lfloor\frac{b_n}{2}\right\rfloor.
            \end{align*}
            床関数の合成の性質と帰納法により,
            \[
              b_n=\left\lfloor\frac{M-6}{2^{n-1}}\right\rfloor.
            \]
            \,$M\ge6$では,十分大きな$n$で中身が$[0,1)$に入り,床は0になる.
            \,$M<6$では,十分大きな$n$で中身が$(-1,0)$に入り,床は$-1$になる.
            これより【解答】の値に有限回で到達する.

            実際,整数$x$が不動点となる条件は,
            \[
              x\le\frac{x+6}{2}<x+1
              \iff 4<x\le6
            \]
            なので,不動点は5と6の二つである.
            初項によって到達する不動点が分かれることは,この後の\sectionref{節：初項依存型}ともつながる.
    \end{enumerate}

